\section{Problem 1.}
\begin{enumerate}[(a)]
\item Let \( a = ( a_1, a_2, ..., a_k ) \in \FF_q^k \) be a vector, \( f : \FF_q^k \to \FF_q \) be a function defined as:
\[
    f_a ( x_1, ..., x_k ) = \sum_{ i = 1 }^k a_i x_i
\]
Show that each element of \( \FF_q \) is the image of exactly \( q^{ k - 1 } \) vectors in \( \FF_q^k \) under the mapping of \( f \).
\begin{solution}
Denote the preimage of an arbitrary vector \( b \in \FF_q^k \) as:
\[
    \inv{f}( v ) :=
    \pair{ x_i \in \FF_q
    \mid \sum_{ i = 1 }^k a_i x_i = v }
\]
Assume the number of solution with fixing the first \( k \) coordinates is \( S_k \), the total number of solution is \( S_0 \), we have the iteration equation:
\[
    S_0 = q S_1
\]
As the same, we have:
\[
    S_i = q S_{ i + 1 }
\]
until \( S_k = 1 \). As a result, the number of solution is exactly:
\[
    S_0 = \underbrace{ q \cdot q \cdot \cdots \cdot q }_{ k + 1 } \cdot 1 = q^{ k - 1 }
\]

\end{solution}

\item Let \( C \) be a \( ( n, k, d ) \) code over \( \FF_q \) and let \( T \) be a \( q^k \times n \) matrix whose rows are codewords of \( C \). Show that each nonzero elements of \( \FF_q \) appears in every nonzero column in \( T \) exactly \( q^{ k - 1 } \) times.
\begin{proof}
    Assume \( x \in \FF_q \) be an element, we prove it by contradiction. If there exist some column who has the number of element \( x \) less than \( q^{ k - 1 } \), as there are totally \( q^{ k } \) rows, there must exist some element \( x' \in \FF_q \) such that it occurs at more than \( q^{k - 1} \) rows. As the rows are codewords of \( C \), we know that such rows containing \( x' \) should be different. As there are \( k - 1 \) coordinates and \( q \) elements for each position, there are at most \( q^{ k - 1 } \) different vectors. As a result, we know that there must be some repeating codewords, which is contradict to the hypothesis. By the contradiction, we get the result that there number cannot be strict less or greater than \( q^{ k - 1 } \).
\end{proof}

\end{enumerate}

\section{Uniformly Random Parity-Check Matrix}
Let \( R \in ( 0, 1 ) \), \( k = \floor{ R n } \), and \( r = n - k \). Choose a matrix \( H \in \FF_q^{ r \times n } \) uniformly at random, and define \( C \) to be the code having \( H \) as its parity-check matrix, i.e.:
\[
    C := \pair{
        x \in \FF_q^n 
        \mid 
        H x^T = 0
    }
\]
Prove that for every fixed \( \delta \) satisfying that \( H( \delta ) < 1 - R \), with high probability \( H \) has rank \( r \) and:
\[
    d( C ) > \delta n
\]
\begin{proof}
We have:

\begin{enumerate}
\item If \( H \) is full rank, we have the facts that the first row is nonzero, the \( k + 1 \) row is not the linear combination of the first \( k \) rows. The probability of the first row to be nonzero is exactly:
\[
    1 - \left( \frac{1}{q} \right)^n
\]
The probability of the second row is not a linear combination (multiplier) of the first row is:
\[
    1 - ( \frac{q}{ q^n } )
\]
As a result, the probability should be:
\[
    P[ H \text{ is full row rank } ]
    = \left( 1 - \left( \frac{1}{q} \right)^n \right) 
    \cdot \left( 1 - \left( \frac{q}{q^n} \right) \right) 
    \cdot \cdots 
    \cdot (1 - \left( \frac{ q^{ r - 1 } }{q} \right))
\]
It's:
\[
    P[ H \text{ is full row rank } ] 
    = \prod_{ i = 0 }^{r  - 1} ( 1 - \frac{ 1 }{ q^{ n - i } } )
\]
which is really high if \( q \) is large.

\item Firstly, we know that:
\[
    d({\cal C}) = \min_{ x \in C } w(x)
\]
Thus, we have:
\[
    P[ d(C) > \delta n ] = 1 - P[ d(C) \le \delta n ] 
\]
We try to check the bound of the right hand side. Considering that:
\[
    P[ d(C) \le \delta n ] 
    = P[ x \in \FF_q^n, x \in C, w(x) \le \delta n ]
\]
As we know that every vector \( x \) in \( C \) satisfy that \( H x^T = 0 \), we get the probability to be:
\[
    P[ x \in \FF_q^n, w(x) \le \delta n, H x^T = 0 ]
\]
For a given row, we have:
\[
    h_i x^T = 0 \iff \sum_{ j = 1 }^n h_{ij} x_j = 0
\]
According to the question 1, we know there are \( q^{ n - 1 } \) solutions. Thus the probability is
\[
    \frac{ q^{ n - 1 } }{ q^{n} } = \frac{ 1 }{ q }
\]
As the rows are independent, the probability of \( H x^T = 0 \) is exactly \( \frac{1}{q^r} \). Thus, we have:
\[
    P[ x \in \FF_q^n, w(x) \le \delta n, H x^T = 0 ] 
    \le \sum_{ x, w(x) \le \delta n } P[ H x^T = 0 ] = 
    \abs{\pair{ x \mid w(x) \le \delta n }} 
    \cdot \frac{1}{q^r}
\]
According to the Hamming bound, we know that:
\[
    \abs{ B( 0, \floor{ \delta n } ) } 
    \le q^{ H_q( \delta ) n }
\]
and \( H_q(\delta) < 1 - R \); thus, we have \( r = n - k = n - \floor{ nR } \ge n( 1 - R ) \):
\[
    P[ d(C) \le \delta n ] \le 
    q^{ H_q(\delta)n - n( 1 - R )  }
    \le q^{ n(1 - R + \varepsilon) - n(1 - R) }
    = q^{ n\varepsilon } \to 0
\]
As a result, we have:
\[
    \lim_{ n \to \infty} 
    P[ d(C) > \delta n ] 
    = 1 - \lim_{ n \to \infty } q^{ n\varepsilon } 
    = 1
\]

\end{enumerate}

\end{proof}



\section{Minimum distance of a code}
Let \( R \in ( 0, 1/2 ) \) and \( M = \ceil{ q^{ Rn } } \), a random q-ary \( ( n, M ) \) code over \( \Sigma \) is obtained by choosing codewords \( X_1, X_2, ..., X_M \in \Sigma^n \) independently, uniformly, and randomly. Let \( \delta(C) \) be the relative minimum distance of the random code, we write \( \inv{ H_q } \) for the inverse of \( H_q \) on the interval \( [ 0, 1 - 1/q ] \).
\begin{enumerate}[(a)]
\item Show that for any fixed \( \varepsilon > 0 \) with \( 1 - 2 R - \varepsilon > 0 \), with high probability we have:
\[
    \delta(C) \ge \inv{ H_q } ( 1 - 2 R - \varepsilon ) - o (1)
\]
\begin{proof}
    Denote \( \delta' = \inv{ H_q }( 1 - 2R - \varepsilon ) \), we firstly show that: 
    \[
        P[ d(C) \le n \delta'] \to 0
    \]
    The probability of \( d(C) \le n\delta' \) is the probability:
    \[
        P [ \exists c_i, c_j, c_i \neq c_j, d( c_i, c_j ) \le n\delta' ]
        = P[ \bigcup_{ i, j } d( c_i, c_j ) \le n\delta' ]
        \le \sum_{ i, j } P[ d( c_i, c_j ) \le n\delta' ]
    \]
    For arbitrary two vectors in the space, we have the probability:
    \[
        P[ d( c_i, c_j ) \le n\delta' ]
        = \frac{ \Vol_q ( n, \floor{ n, n \delta' } ) }{ q^n }
        \le q^{ n H_q( \delta' ) - o(n) - n }
    \]
    Thus, the probability is:
    \[
        P[ d(c) \le n \delta' ] 
        \le \begin{pmatrix}
            M \\ 
            2 \\
        \end{pmatrix} \cdot 
        q^{ n( 1 - 2R - \varepsilon ) }
        = \frac{ M ( M - 1 ) }{ 2 }
        \cdot q^{ n ( 1 - 2R - \varepsilon ) - o(n) - n }
        \le q^{ 2Rn + n( 1 - 2R - \varepsilon ) - o(n) - n}
    \]
    As \( \varepsilon \) satisfies the equation \( 1 - 2R - \varepsilon > 0 \), we have:
    \[
        P[ d( C ) \le n \delta' ] 
        \le q^{ -\varepsilon n - o(n) } \to 0 
        \quad 
        ( n \to \infty )
    \]
    As a result, we have high probability that;
    \[
        d(C) \ge n\inv{ H_q }( 1 - 2R - \varepsilon )
    \]
    which implies:
    \[
        \delta(C) 
        \ge \inv{H_q}( 1 - 2 R - \varepsilon )
        \ge \inv{H_q}( 1 - 2 R - \varepsilon )
        - o(1)
    \]
\end{proof}


\item Show that for any fixed \( \varepsilon > 0 \), with high probability we have:
\[
    \delta(C) 
    \le \inv{ H_q } ( 1 - 2R ) + \varepsilon
\]
\begin{proof}
    It's equivalent to check:
    \[
        P[ d(C) \le n \delta' ]
    \]
    where \( \delta' = \inv{ H_q }(1 - 2R) \). Considering that there are \( C_M^2 \) pairs of distance and each pair has probability:
    \[
        P[ d( c_i, c_j ) \le n\delta' ] 
        \le q^{ nH_q(\delta') - n -o(n) }
        \le q^{ n(1 - 2R + H_q( \varepsilon )) - n - o(n) }
    \]
    According to the (a), we have the probability:
    \[ 
        P[ d(C) \le n\delta' ] 
        \le \begin{pmatrix}
            M \\ 
            2 \\
        \end{pmatrix}
        q^{ nH_q( \varepsilon ) -2nR }
        \le 
        \frac{1}{2} q^{ 2nR } q^{ nH_q( \varepsilon ) - 2nR }
        \le 
        q^{ nH_q( \varepsilon ) }
    \]
    As \( H_q( \varepsilon ) \) is a positive number, we know the bound tends to infinity with \( n \), which means we have high probability the inequality holds:
    \[
        \delta(C) \le \inv{ H_q }( 1 - 2R ) + \varepsilon
    \]
\end{proof}





\end{enumerate}

\section{Hamming Ball}
Denote the volume of a Hamming Ball of radius in \( \FF_q^n \) by:
\[
    V_q( n, t ) := \sum_{ i = 0 }^t \begin{pmatrix}
        n \\ 
        i \\
    \end{pmatrix} 
    \cdot 
    ( q - 1 )^i
\]
Let \( r = n - k \), and let \( D \) be an integer satisfying:
\[
    V_q( n - 1, D - 2 ) < q^r 
\]
Prove that one can construct a \( r \times n \) parity-check matrix \( H \) of an \( ( n, k, d ) \) linear code over \( \FF_q \) with \( d \ge D \).

\begin{proof}
To provide the minimum distance \( d \ge D \); we prove for arbitrary set of less equal than \( D - 1 \) columns, it's linearly independent. To do this, we can choose the columns linearly independently; but it need to require there are enough vectors to choose. As we only require \( d \ge D \), we need at least \( D - 1 \) linearly independent columns. When we have already choose \( k \) columns, the vector has the form:
\[
    v = c_1 v_1 + c_2 v_2 + \cdots c_k v_k, c_i \in \FF_q
\]
cannot be chosen. As we only need to check until the \( D - 1 \) column, we know that there are at most \( D - 2 \) coefficients. When it has \( D - 2 \) coefficients, the most difficult case, we have the linearly dependent set to be bounded by:
\[
    \pair{ v \in \FF_q^{ k - 1 } \mid V_q( k - 1, D - 2 ) }
\]
Use the loss bound, we know that if it is strict less than the whole space, we can choose a new. We check the last column:
\[
    V_q( n - 1, D - 2 ) < q^r
\]
holds by the hypothesis. As a result, we know there are new vector we can choose, which means we finally find \( n \) vectors.


\end{proof}




\section{Q5}
Let \( \alpha = ( \alpha_1, \alpha_2, ..., \alpha_n ) \in \FF_q^n \) be a vector with distinct entries, and \( v \in (F^*_q)^n \), prove that for any \( 1 \le k \le n - 1 \), the dual code of \( \text{GRS}_q( \alpha, k, v ) \) is \( \text{GRS}_q( \alpha, n - k, v' ) \) for some \( v' = ( v_1', v_2', ..., v_n' ) \in (\FF_q^*)^n \) 

\begin{proof}
Considering that the generator matrix of \( GRS_q( \alpha, k, v ) \) code is:
\[
    \begin{pmatrix}
        v_1 & v_2 & \cdots & v_n \\
        v_1 \alpha_1 & v_2 \alpha_2 & \cdots & v_n \alpha_n \\
        \vdots & & & \vdots \\
        v_1 \alpha_1^{ k - 1 } & v_2 \alpha_2^{ k - 1 } & \cdots & v_n \alpha_n^{ k - 1 } \\
    \end{pmatrix}
\]
Assume the \( H \) exist, it should has the generator matrix:
\[
    \begin{pmatrix}
        v'_1 & v'_2 & \cdots & v'_n \\
        v'_1 \alpha_1 & v'_2 \alpha_2 & \cdots & v'_n \alpha_n \\
        \vdots & & & \vdots \\
        v'_1 \alpha_1^{ n - k - 1 } & v'_2 \alpha_2^{ n - k - 1 } & \cdots & v'_n \alpha_n^{ n - k - 1 } \\
    \end{pmatrix}
\]
To make them dual, we need \( GH^T = 0 \), which means:
\[
    ( GH^T )_{ ij } = \sum_{ k = 1 }^n v_k v'_k a_k^{ i } a_k^{ j } = \sum_{ k = 1 }^n v_k v_k' a_k^{ i + j } = 0
\]
Thus we have:
\[
    \sum_{ k = 1 }^n ( v_k v_k' ) a_k^{ l } = 0, l \in \pair{ 0, ..., n - 2 }
\]
We need the \( ( v_k v_k' ) \) to be a solution for the \( ( n - 1 ) \times n \) Vandermonde matrix:
\[
    \begin{pmatrix}
    1 & 1 & \dots & 1 \\
    \alpha_1 & \alpha_2 & \dots & \alpha_n \\
    \vdots & \vdots & \ddots & \vdots \\ 
    \alpha_1^{n-2} & \alpha_2^{n-2} & \dots & \alpha_n^{n-2} \\
    \end{pmatrix}
\]
According to the knowledge of linear algebra, we know that:
\[
    v_k v_k' = \frac{ 1 }{ \left( \prod_{ i = 1 }^n ( x - \alpha_i ) \right)'( \alpha_k ) }
\]
is a solution. In above equation, we consider the \( x \) to be the variable of a polynomial and we evaluate at \( a_k \) after derivative. In conclusion, the \( H \) exists and \( v_k' \) has the expression:
\[
    v_k' = \frac{ 1 }{ v_k P'( \alpha_k ) },
    \quad 
    P(x) = \prod_{ i = 1 }^n ( x - \alpha_i )
\]



\end{proof}


